\section{Syntax}

% BEGIN GENERATED ARTIFACT: def:truth-value-propositional-logic
\begin{tcolorbox}[colback=propbox, colframe=propborder, arc=2pt,
  left=6pt, right=6pt, top=4pt, bottom=4pt,
  title={\small\textbf{Definition (Truth Value)}},
  fonttitle=\small\bfseries]
\begin{definition}[Truth Value]
\label{def:truth-value-propositional-logic}
The set of truth values is
\[
\mathbb{B} := \{\mathrm{False},\mathrm{True}\}.
\]
Some sources identify this set with \(\{0,1\}\), with \(0\) read as false and \(1\) read as true.
\end{definition}
\end{tcolorbox}

\begin{remark*}[Standard quantified statement]
\[
\mathbb{B}=\{\mathrm{False},\mathrm{True}\}.
\]
\end{remark*}

\begin{remark*}[Interpretation]
Truth values are the semantic targets of propositional formulas. Syntax determines which strings are formulas; semantics later assigns elements of \(\mathbb{B}\) to formulas.
\end{remark*}

\begin{dependencies}
\NoLocalDependencies
\end{dependencies}
% END GENERATED ARTIFACT: def:truth-value-propositional-logic
% BEGIN GENERATED ARTIFACT: def:propositional-language
\begin{tcolorbox}[colback=propbox, colframe=propborder, arc=2pt,
  left=6pt, right=6pt, top=4pt, bottom=4pt,
  title={\small\textbf{Definition (Propositional Language)}},
  fonttitle=\small\bfseries]
\begin{definition}[Propositional Language]
\label{def:propositional-language}
A propositional language consists of:
\begin{enumerate}
\item a set \(\Prop\) of propositional variables;
\item logical connectives;
\item punctuation symbols such as parentheses;
\item recursive formation rules specifying the well-formed formulas.
\end{enumerate}
\end{definition}
\end{tcolorbox}

\begin{remark*}[Standard quantified statement]
\[
\mathcal L_{\mathrm{prop}}
=
(\Prop,\{\neg,\land,\lor,\to,\leftrightarrow\},\text{punctuation},\text{formation rules}).
\]
\end{remark*}

\begin{remark*}[Interpretation]
The language supplies symbols and grammar. It does not yet assign truth values, truth conditions, or logical consequence.
\end{remark*}

\begin{dependencies}
\begin{itemize}
  \item \hyperref[def:truth-value-propositional-logic]{Truth Value}
\end{itemize}
\end{dependencies}
% END GENERATED ARTIFACT: def:propositional-language
% BEGIN GENERATED ARTIFACT: def:propositional-variable
\begin{tcolorbox}[colback=propbox, colframe=propborder, arc=2pt,
  left=6pt, right=6pt, top=4pt, bottom=4pt,
  title={\small\textbf{Definition (Propositional Variable)}},
  fonttitle=\small\bfseries]
\begin{definition}[Propositional Variable]
\label{def:propositional-variable}
A propositional variable is an atomic symbol intended to stand for a proposition. A standard supply of variables is
\[
\Prop=\{P_1,P_2,P_3,\ldots\}.
\]
Informally, variables are often written \(P,Q,R,S,\ldots\).
\end{definition}
\end{tcolorbox}

\begin{remark*}[Standard quantified statement]
\[
\operatorname{PropositionalVariable}(P)\Longleftrightarrow P\in\Prop.
\]
\end{remark*}

\begin{remark*}[Definition predicate reading]
\(\operatorname{PropositionalVariable}(P)\) means \(P\in\Prop\).
\end{remark*}

\begin{remark*}[Interpretation]
A propositional variable has no internal syntactic structure in propositional logic. It is a basic symbol from which more complex formulas are formed.
\end{remark*}

\begin{remark*}[Examples]
\(P,Q,R,P_1,P_2\) are propositional variables when they belong to \(\Prop\).
\end{remark*}

\begin{remark*}[Non-examples]
\(\neg P\), \(P\land Q\), and \((P\to Q)\) are not propositional variables; they are compound formulas.
\end{remark*}

\begin{dependencies}
\begin{itemize}
  \item \hyperref[def:propositional-language]{Propositional Language}
\end{itemize}
\end{dependencies}
% END GENERATED ARTIFACT: def:propositional-variable
% BEGIN GENERATED ARTIFACT: def:logical-connective
\begin{tcolorbox}[colback=propbox, colframe=propborder, arc=2pt,
  left=6pt, right=6pt, top=4pt, bottom=4pt,
  title={\small\textbf{Definition (Logical Connective)}},
  fonttitle=\small\bfseries]
\begin{definition}[Logical Connective]
\label{def:logical-connective}
The standard primitive propositional connectives are
\[
\neg,\qquad \land,\qquad \lor,\qquad \to,\qquad \leftrightarrow.
\]
The connective \(\neg\) is unary. The connectives \(\land,\lor,\to,\leftrightarrow\) are binary.

\[
\begin{array}{c|c|c|c|c}
\text{Symbol} & \text{Name} & \text{Arity} & \text{Formation pattern} & \text{Reading}\\
\hline
\neg & \text{negation} & \text{unary} & \neg\varphi & \text{not }\varphi\\
\land & \text{conjunction} & \text{binary} & (\varphi\land\psi) & \varphi\text{ and }\psi\\
\lor & \text{disjunction} & \text{binary} & (\varphi\lor\psi) & \varphi\text{ or }\psi\\
\to & \text{conditional} & \text{binary} & (\varphi\to\psi) & \text{if }\varphi,\text{ then }\psi\\
\leftrightarrow & \text{biconditional} & \text{binary} & (\varphi\leftrightarrow\psi) & \varphi\text{ iff }\psi
\end{array}
\]
\end{definition}
\end{tcolorbox}

\begin{remark*}[Standard quantified statement]
\[
\operatorname{UnaryConnective}(\neg)
\quad\text{and}\quad
\operatorname{BinaryConnective}(\land)\land \operatorname{BinaryConnective}(\lor)\land \operatorname{BinaryConnective}(\to)\land \operatorname{BinaryConnective}(\leftrightarrow).
\]
\end{remark*}

\begin{remark*}[Interpretation]
Arity records how many formulas a connective uses to form a new formula. The truth behavior of these connectives belongs to semantics, not to the syntactic definition.
\end{remark*}

\begin{dependencies}
\begin{itemize}
  \item \hyperref[def:propositional-language]{Propositional Language}
\end{itemize}
\end{dependencies}
% END GENERATED ARTIFACT: def:logical-connective
% BEGIN GENERATED ARTIFACT: def:well-formed-formula
\begin{tcolorbox}[colback=propbox, colframe=propborder, arc=2pt,
  left=6pt, right=6pt, top=4pt, bottom=4pt,
  title={\small\textbf{Definition (Well-Formed Formula)}},
  fonttitle=\small\bfseries]
\begin{definition}[Well-Formed Formula]
\label{def:well-formed-formula}
The set \(\WFF\) of well-formed formulas is the smallest set of strings satisfying the following formation rules:
\begin{enumerate}
\item Every propositional variable \(P\in\Prop\) belongs to \(\WFF\).
\item If \(\varphi\in\WFF\), then \(\neg\varphi\in\WFF\).
\item If \(\varphi,\psi\in\WFF\) and \(\circ\in\{\land,\lor,\to,\leftrightarrow\}\), then \((\varphi\circ\psi)\in\WFF\).
\item No string is a well-formed formula unless it is obtained by finitely many applications of the previous clauses.
\end{enumerate}
\end{definition}
\end{tcolorbox}

\begin{remark*}[Standard quantified statement]
\[
\WFF
=
\bigcap
\left\{
S :
\Prop\subseteq S,\;
(\forall\varphi\in S)(\neg\varphi\in S),\;
(\forall\varphi,\psi\in S)
(\forall\circ\in\{\land,\lor,\to,\leftrightarrow\})
\bigl((\varphi\circ\psi)\in S\bigr)
\right\}.
\]
\end{remark*}

\begin{remark*}[Definition predicate reading]
\(\operatorname{WellFormedFormula}(\varphi)\) means \(\varphi\in\WFF\), read as “\(\varphi\) is a well-formed formula.”
\end{remark*}

\begin{remark*}[Negated quantified statement]
\[
\sigma\notin\WFF
\Longleftrightarrow
\exists S
\left[
\Prop\subseteq S
\land
(\forall\varphi\in S)(\neg\varphi\in S)
\land
(\forall\varphi,\psi\in S)
(\forall\circ\in\{\land,\lor,\to,\leftrightarrow\})
((\varphi\circ\psi)\in S)
\land
\sigma\notin S
\right].
\]
\end{remark*}

\begin{remark*}[Negation predicate reading]
\(\neg \operatorname{WellFormedFormula}(\sigma)\) means the string \(\sigma\) is excluded by at least one formation-closed set containing \(\Prop\); equivalently, it is not generated by finitely many applications of the formation rules.
\end{remark*}

\begin{remark*}[Failure mode decomposition]
\[
\underbrace{P\land}_{\text{missing right operand}}
\qquad
\underbrace{(P\ Q)}_{\text{missing binary connective}}
\qquad
\underbrace{\neg\lor P}_{\text{negation applied to a non-formula}}.
\]
\end{remark*}

\begin{remark*}[Interpretation]
Well-formed formulas are exactly the strings to which the semantic and proof-theoretic machinery of propositional logic applies.
\end{remark*}

\begin{remark*}[Examples]
\(P\), \(\neg P\), \((P\to Q)\), and \((P\to(Q\lor R))\) are well-formed formulas.
\end{remark*}

\begin{remark*}[Non-examples]
\(P\land\), \((P Q)\), and \(\neg\lor P\) are not well-formed formulas.
\end{remark*}

\begin{dependencies}
\begin{itemize}
  \item \hyperref[def:propositional-variable]{Propositional Variable}
  \item \hyperref[def:logical-connective]{Logical Connective}
\end{itemize}
\end{dependencies}
% END GENERATED ARTIFACT: def:well-formed-formula
% BEGIN GENERATED ARTIFACT: lem:closure-under-formation-propositional-formulas
\begin{tcolorbox}[colback=lembox, colframe=lemborder, arc=2pt,
  left=6pt, right=6pt, top=4pt, bottom=4pt,
  title={\small\textbf{Lemma (Closure Under Formation)}},
  fonttitle=\small\bfseries]
\begin{lemma}[Closure Under Formation]
\label{lem:closure-under-formation-propositional-formulas}
The set \(\WFF\) is closed under the propositional formation rules.

That is, if \(\varphi\in\WFF\), then
\[
\neg\varphi\in\WFF.
\]
If \(\varphi,\psi\in\WFF\) and
\[
\circ\in\{\land,\lor,\to,\leftrightarrow\},
\]
then
\[
(\varphi\circ\psi)\in\WFF.
\]

\noindent\hyperref[prf:closure-under-formation-propositional-formulas]{Go to proof.}
\end{lemma}
\end{tcolorbox}

\begin{remark*}[Standard quantified statement]
\[
(\forall\varphi\in\WFF)(\neg\varphi\in\WFF)
\]
and
\[
(\forall\varphi,\psi\in\WFF)
(\forall\circ\in\{\land,\lor,\to,\leftrightarrow\})
((\varphi\circ\psi)\in\WFF).
\]
\end{remark*}

\begin{remark*}[Interpretation]
Closure says that every legal use of a formation rule produces another
well-formed formula. It is the forward direction of the recursive grammar.
\end{remark*}

\begin{dependencies}
\begin{itemize}
  \item \hyperref[def:well-formed-formula]{Well-Formed Formula}
\end{itemize}
\end{dependencies}
% END GENERATED ARTIFACT: lem:closure-under-formation-propositional-formulas
% BEGIN GENERATED ARTIFACT: thm:minimality-of-well-formed-formulas
\begin{tcolorbox}[colback=thmbox, colframe=thmborder, arc=2pt,
  left=6pt, right=6pt, top=4pt, bottom=4pt,
  title={\small\textbf{Theorem (Minimality of Well-Formed Formulas)}},
  fonttitle=\small\bfseries]
\begin{theorem}[Minimality of Well-Formed Formulas]
\label{thm:minimality-of-well-formed-formulas}
Let \(S\) be a set of strings over the propositional alphabet. Suppose:
\begin{enumerate}
  \item \(\Prop\subseteq S\);
  \item if \(\varphi\in S\), then \(\neg\varphi\in S\);
  \item if \(\varphi,\psi\in S\) and
        \(\circ\in\{\land,\lor,\to,\leftrightarrow\}\), then
        \((\varphi\circ\psi)\in S\).
\end{enumerate}
Then
\[
\WFF\subseteq S.
\]

\noindent\hyperref[prf:minimality-of-well-formed-formulas]{Go to proof.}
\end{theorem}
\end{tcolorbox}

\begin{remark*}[Standard quantified statement]
\[
\begin{aligned}
(\forall S)\Big[
&\Prop\subseteq S
\land
(\forall\varphi\in S)(\neg\varphi\in S)\\
&\land
(\forall\varphi,\psi\in S)
(\forall\circ\in\{\land,\lor,\to,\leftrightarrow\})
((\varphi\circ\psi)\in S)
\Big]
\to
\WFF\subseteq S.
\end{aligned}
\]
\end{remark*}

\begin{remark*}[Interpretation]
Minimality is the formal content of the clause ``nothing else is a wff.'' It says
that \(\WFF\) is contained in every set of strings that contains the
propositional variables and is closed under the formation rules.
\end{remark*}

\begin{dependencies}
\begin{itemize}
  \item \hyperref[def:well-formed-formula]{Well-Formed Formula}
\end{itemize}
\end{dependencies}
% END GENERATED ARTIFACT: thm:minimality-of-well-formed-formulas
% BEGIN GENERATED ARTIFACT: def:atomic-formula-propositional-logic
\begin{tcolorbox}[colback=propbox, colframe=propborder, arc=2pt,
  left=6pt, right=6pt, top=4pt, bottom=4pt,
  title={\small\textbf{Definition (Atomic Formula)}},
  fonttitle=\small\bfseries]
\begin{definition}[Atomic Formula]
\label{def:atomic-formula-propositional-logic}
An atomic formula is a well-formed formula that is a propositional variable.
\end{definition}
\end{tcolorbox}

\begin{remark*}[Standard quantified statement]
\[
\operatorname{AtomicFormula}(\varphi)
\Longleftrightarrow
\varphi\in\WFF\land\varphi\in\Prop.
\]
\end{remark*}

\begin{remark*}[Definition predicate reading]
\(\operatorname{AtomicFormula}(\varphi)\) means that \(\varphi\) is a propositional variable viewed as a well-formed formula.
\end{remark*}

\begin{remark*}[Interpretation]
Atomic formulas are the base cases of the recursive grammar. They contain no connective structure.
\end{remark*}

\begin{dependencies}
\begin{itemize}
  \item \hyperref[def:well-formed-formula]{Well-Formed Formula}
  \item \hyperref[def:propositional-variable]{Propositional Variable}
\end{itemize}
\end{dependencies}
% END GENERATED ARTIFACT: def:atomic-formula-propositional-logic
% BEGIN GENERATED ARTIFACT: def:molecular-formula-propositional-logic
\begin{tcolorbox}[colback=propbox, colframe=propborder, arc=2pt,
  left=6pt, right=6pt, top=4pt, bottom=4pt,
  title={\small\textbf{Definition (Molecular Formula)}},
  fonttitle=\small\bfseries]
\begin{definition}[Molecular Formula]
\label{def:molecular-formula-propositional-logic}
A molecular formula is a well-formed formula that is not atomic. Equivalently, it is a well-formed formula formed using at least one logical connective.
\end{definition}
\end{tcolorbox}

\begin{remark*}[Standard quantified statement]
\[
\operatorname{MolecularFormula}(\varphi)
\Longleftrightarrow
\varphi\in\WFF\land\neg \operatorname{AtomicFormula}(\varphi).
\]
\end{remark*}

\begin{remark*}[Definition predicate reading]
\(\operatorname{MolecularFormula}(\varphi)\) means that \(\varphi\) has nontrivial connective structure.
\end{remark*}

\begin{remark*}[Interpretation]
Molecular formulas are exactly the formulas obtained by applying negation or a binary connective at least once.
\end{remark*}

\begin{dependencies}
\begin{itemize}
  \item \hyperref[def:well-formed-formula]{Well-Formed Formula}
  \item \hyperref[def:atomic-formula-propositional-logic]{Atomic Formula}
  \item \hyperref[def:logical-connective]{Logical Connective}
\end{itemize}
\end{dependencies}
% END GENERATED ARTIFACT: def:molecular-formula-propositional-logic
% BEGIN GENERATED ARTIFACT: def:main-connective-propositional-logic
\begin{tcolorbox}[colback=propbox, colframe=propborder, arc=2pt,
  left=6pt, right=6pt, top=4pt, bottom=4pt,
  title={\small\textbf{Definition (Main Connective)}},
  fonttitle=\small\bfseries]
\begin{definition}[Main Connective]
\label{def:main-connective-propositional-logic}
The main connective of a molecular formula is the connective introduced at the final formation step:
\begin{enumerate}
\item for \(\neg\varphi\), the main connective is \(\neg\);
\item for \((\varphi\circ\psi)\), the main connective is \(\circ\).
\end{enumerate}
Atomic formulas have no main connective.
\end{definition}
\end{tcolorbox}

\begin{remark*}[Standard quantified statement]
\[
\operatorname{MainConnective}(\neg\varphi,\neg)
\quad\text{and}\quad
\operatorname{MainConnective}((\varphi\circ\psi),\circ).
\]
\end{remark*}

\begin{remark*}[Definition predicate reading]
\(\operatorname{MainConnective}(\varphi,c)\) means that \(c\) is the outermost connective introduced in the final formation step of \(\varphi\).
\end{remark*}

\begin{remark*}[Interpretation]
The main connective identifies the outermost syntactic operation. It is the first connective used when analyzing a formula from the outside inward.
\end{remark*}

\begin{dependencies}
\begin{itemize}
  \item \hyperref[def:well-formed-formula]{Well-Formed Formula}
  \item \hyperref[def:molecular-formula-propositional-logic]{Molecular Formula}
  \item \hyperref[def:logical-connective]{Logical Connective}
\end{itemize}
\end{dependencies}
% END GENERATED ARTIFACT: def:main-connective-propositional-logic
% BEGIN GENERATED ARTIFACT: thm:structural-induction-propositional-formulas
\begin{tcolorbox}[colback=thmbox, colframe=thmborder, arc=2pt,
  left=6pt, right=6pt, top=4pt, bottom=4pt,
  title={\small\textbf{Theorem (Structural Induction for Propositional Formulas)}},
  fonttitle=\small\bfseries]
\begin{theorem}[Structural Induction for Propositional Formulas]
\label{thm:structural-induction-propositional-formulas}
Let \(A(\varphi)\) be a property of well-formed formulas. Suppose:
\begin{enumerate}
\item \(A(P)\) holds for every \(P\in\Prop\);
\item if \(A(\varphi)\) holds, then \(A(\neg\varphi)\) holds;
\item if \(A(\varphi)\) and \(A(\psi)\) hold, then \(A((\varphi\circ\psi))\) holds for every \(\circ\in\{\land,\lor,\to,\leftrightarrow\}\).
\end{enumerate}
Then \(A(\theta)\) holds for every \(\theta\in\WFF\).

\noindent\hyperref[prf:structural-induction-propositional-formulas]{Go to proof.}
\end{theorem}
\end{tcolorbox}

\begin{remark*}[Standard quantified statement]
\[
\left[
(\forall P\in\Prop)A(P)
\land
(\forall\varphi\in\WFF)(A(\varphi)\to A(\neg\varphi))
\land
(\forall\varphi,\psi\in\WFF)
(\forall\circ\in\{\land,\lor,\to,\leftrightarrow\})
\bigl((A(\varphi)\land A(\psi))\to A((\varphi\circ\psi))\bigr)
\right]
\to
(\forall\theta\in\WFF)A(\theta).
\]
\end{remark*}

\begin{remark*}[Interpretation]
To prove a statement for all formulas, it suffices to prove it for propositional variables and show it is preserved by each formula-forming operation.
\end{remark*}

\begin{dependencies}
\begin{itemize}
  \item \hyperref[def:well-formed-formula]{Well-Formed Formula}
\end{itemize}
\end{dependencies}
% END GENERATED ARTIFACT: thm:structural-induction-propositional-formulas
% BEGIN GENERATED ARTIFACT: thm:structural-recursion-propositional-formulas
\begin{tcolorbox}[colback=thmbox, colframe=thmborder, arc=2pt,
  left=6pt, right=6pt, top=4pt, bottom=4pt,
  title={\small\textbf{Theorem (Structural Recursion for Propositional Formulas)}},
  fonttitle=\small\bfseries]
\begin{theorem}[Structural Recursion for Propositional Formulas]
\label{thm:structural-recursion-propositional-formulas}
To define a function \(F\) on \(\WFF\), it suffices to specify:
\begin{enumerate}
\item \(F(P)\) for every \(P\in\Prop\);
\item \(F(\neg\varphi)\) in terms of \(F(\varphi)\);
\item \(F((\varphi\circ\psi))\) in terms of \(F(\varphi)\), \(F(\psi)\), and \(\circ\).
\end{enumerate}
When these clauses respect the formation rules, they determine a unique function on \(\WFF\).

\noindent\hyperref[prf:structural-recursion-propositional-formulas]{Go to proof.}
\end{theorem}
\end{tcolorbox}

\begin{remark*}[Standard quantified statement]
\[
\text{recursive data on the formation clauses of }\WFF
\Longrightarrow
\exists!F:\WFF\to X
\text{ satisfying those clauses}.
\]
\end{remark*}

\begin{remark*}[Interpretation]
Structural recursion justifies recursive definitions such as variables of a formula, subformulas, depth, connective count, syntax tree, and truth evaluation.
\end{remark*}

\begin{dependencies}
\begin{itemize}
  \item \hyperref[def:well-formed-formula]{Well-Formed Formula}
  \item \hyperref[thm:structural-induction-propositional-formulas]{Structural Induction for Propositional Formulas}
\end{itemize}
\end{dependencies}
% END GENERATED ARTIFACT: thm:structural-recursion-propositional-formulas
% BEGIN GENERATED ARTIFACT: lem:constructor-disjointness-propositional-formulas
\begin{tcolorbox}[colback=lembox, colframe=lemborder, arc=2pt,
  left=6pt, right=6pt, top=4pt, bottom=4pt,
  title={\small\textbf{Lemma (Constructor Disjointness for Propositional Formulas)}},
  fonttitle=\small\bfseries]
\begin{lemma}[Constructor Disjointness for Propositional Formulas]
\label{lem:constructor-disjointness-propositional-formulas}
The three outer formation cases for propositional formulas are disjoint:
\begin{enumerate}
\item no propositional variable is a negation formula;
\item no propositional variable is a binary formula;
\item no negation formula is a binary formula.
\end{enumerate}

\noindent\hyperref[prf:constructor-disjointness-propositional-formulas]{Go to proof.}
\end{lemma}
\end{tcolorbox}

\begin{remark*}[Standard quantified statement]
\[
P\in\Prop\Rightarrow
\left[
(\forall\varphi\in\WFF)(P\ne\neg\varphi)
\land
(\forall\varphi,\psi\in\WFF)(\forall\circ)(P\ne(\varphi\circ\psi))
\right]
\]
and
\[
(\forall\varphi,\psi,\chi\in\WFF)(\forall\circ)
\bigl(\neg\varphi\ne(\psi\circ\chi)\bigr).
\]
\end{remark*}

\begin{remark*}[Readable decomposition]
\[
\underbrace{P}_{\text{atomic}}
\qquad
\underbrace{\neg\varphi}_{\text{negation case}}
\qquad
\underbrace{(\varphi\circ\psi)}_{\text{binary case}}
\]
belong to three distinct outer-form classes.
\end{remark*}

\begin{remark*}[Interpretation]
A formula cannot be parsed as two different outer formation types.
\end{remark*}

\begin{dependencies}
\begin{itemize}
  \item \hyperref[def:well-formed-formula]{Well-Formed Formula}
  \item \hyperref[def:atomic-formula-propositional-logic]{Atomic Formula}
  \item \hyperref[def:main-connective-propositional-logic]{Main Connective}
\end{itemize}
\end{dependencies}
% END GENERATED ARTIFACT: lem:constructor-disjointness-propositional-formulas
% BEGIN GENERATED ARTIFACT: lem:constructor-injectivity-propositional-formulas
\begin{tcolorbox}[colback=lembox, colframe=lemborder, arc=2pt,
  left=6pt, right=6pt, top=4pt, bottom=4pt,
  title={\small\textbf{Lemma (Constructor Injectivity for Propositional Formulas)}},
  fonttitle=\small\bfseries]
\begin{lemma}[Constructor Injectivity for Propositional Formulas]
\label{lem:constructor-injectivity-propositional-formulas}
The propositional formula constructors are injective:
\begin{enumerate}
\item if \(\neg\varphi=\neg\psi\), then \(\varphi=\psi\);
\item if \((\varphi\circ\psi)=(\alpha\star\beta)\), then \(\varphi=\alpha\), \(\psi=\beta\), and \(\circ=\star\).
\end{enumerate}

\noindent\hyperref[prf:constructor-injectivity-propositional-formulas]{Go to proof.}
\end{lemma}
\end{tcolorbox}

\begin{remark*}[Standard quantified statement]
\[
(\forall\varphi,\psi\in\WFF)
(\neg\varphi=\neg\psi\to\varphi=\psi)
\]
and
\[
(\forall\varphi,\psi,\alpha,\beta\in\WFF)
(\forall\circ,\star\in\{\land,\lor,\to,\leftrightarrow\})
\left[
(\varphi\circ\psi)=(\alpha\star\beta)
\to
(\varphi=\alpha\land\psi=\beta\land\circ=\star)
\right].
\]
\end{remark*}

\begin{remark*}[Interpretation]
If two formulas are built by the same outer constructor and are equal as strings, then their immediate constituents are equal.
\end{remark*}

\begin{dependencies}
\begin{itemize}
  \item \hyperref[def:well-formed-formula]{Well-Formed Formula}
  \item \hyperref[def:logical-connective]{Logical Connective}
\end{itemize}
\end{dependencies}
% END GENERATED ARTIFACT: lem:constructor-injectivity-propositional-formulas
% BEGIN GENERATED ARTIFACT: thm:unique-decomposition-propositional-formulas
\begin{tcolorbox}[colback=thmbox, colframe=thmborder, arc=2pt,
  left=6pt, right=6pt, top=4pt, bottom=4pt,
  title={\small\textbf{Theorem (Unique Decomposition for Propositional Formulas)}},
  fonttitle=\small\bfseries]
\begin{theorem}[Unique Decomposition for Propositional Formulas]
\label{thm:unique-decomposition-propositional-formulas}
Every \(\varphi\in\WFF\) has exactly one outermost form:
\begin{enumerate}
\item \(\varphi=P\) for a unique \(P\in\Prop\);
\item \(\varphi=\neg\psi\) for a unique \(\psi\in\WFF\);
\item \(\varphi=(\psi\circ\chi)\) for unique \(\psi,\chi\in\WFF\) and a unique binary connective \(\circ\in\{\land,\lor,\to,\leftrightarrow\}\).
\end{enumerate}

\noindent\hyperref[prf:unique-decomposition-propositional-formulas]{Go to proof.}
\end{theorem}
\end{tcolorbox}

\begin{remark*}[Standard quantified statement]
\[
\forall\varphi\in\WFF,\quad
\text{exactly one of the following holds:}
\]
\[
\varphi\in\Prop,\qquad
\exists!\psi\in\WFF\,(\varphi=\neg\psi),\qquad
\exists!\psi,\chi\in\WFF\;\exists!\circ\in\{\land,\lor,\to,\leftrightarrow\}\,(\varphi=(\psi\circ\chi)).
\]
\end{remark*}

\begin{remark*}[Readable decomposition]
\[
\varphi
\in
\underbrace{\Prop}_{\text{atomic case}}
\quad\text{or}\quad
\varphi=
\underbrace{\neg\psi}_{\text{negation case}}
\quad\text{or}\quad
\varphi=
\underbrace{(\psi\circ\chi)}_{\text{binary case}},
\]
and exactly one of these alternatives holds.
\end{remark*}

\begin{remark*}[Interpretation]
Unique decomposition makes recursive definitions and induction on formulas unambiguous.
\end{remark*}

\begin{dependencies}
\begin{itemize}
  \item \hyperref[def:well-formed-formula]{Well-Formed Formula}
  \item \hyperref[lem:constructor-disjointness-propositional-formulas]{Constructor Disjointness for Propositional Formulas}
  \item \hyperref[lem:constructor-injectivity-propositional-formulas]{Constructor Injectivity for Propositional Formulas}
\end{itemize}
\end{dependencies}
% END GENERATED ARTIFACT: thm:unique-decomposition-propositional-formulas
% BEGIN GENERATED ARTIFACT: def:parse-tree-propositional-logic
\begin{tcolorbox}[colback=propbox, colframe=propborder, arc=2pt,
  left=6pt, right=6pt, top=4pt, bottom=4pt,
  title={\small\textbf{Definition (Parse Tree)}},
  fonttitle=\small\bfseries]
\begin{definition}[Parse Tree]
\label{def:parse-tree-propositional-logic}
The parse tree, or syntax tree, of a well-formed formula is the tree representation of its recursive construction:
\begin{enumerate}
\item a propositional variable is represented by a single leaf;
\item a formula \(\neg\varphi\) is represented by a root labeled \(\neg\) with one child, the parse tree of \(\varphi\);
\item a formula \((\varphi\circ\psi)\) is represented by a root labeled \(\circ\) with two children, the parse trees of \(\varphi\) and \(\psi\).
\end{enumerate}
\end{definition}
\end{tcolorbox}

\begin{remark*}[Standard quantified statement]
\[
\operatorname{Tree}:\WFF\to\{\text{finite rooted labeled trees}\}
\]
is defined recursively by the three formation clauses for propositional formulas.
\end{remark*}

\begin{remark*}[Interpretation]
The parse tree records the same structure described by unique decomposition. It turns the recursive construction of a formula into a visible tree.
\end{remark*}

\begin{dependencies}
\begin{itemize}
  \item \hyperref[def:well-formed-formula]{Well-Formed Formula}
  \item \hyperref[thm:unique-decomposition-propositional-formulas]{Unique Decomposition for Propositional Formulas}
\end{itemize}
\end{dependencies}
% END GENERATED ARTIFACT: def:parse-tree-propositional-logic
% BEGIN GENERATED ARTIFACT: thm:unique-syntax-tree-propositional-formulas
\begin{tcolorbox}[colback=thmbox, colframe=thmborder, arc=2pt,
  left=6pt, right=6pt, top=4pt, bottom=4pt,
  title={\small\textbf{Theorem (Unique Syntax Tree for Propositional Formulas)}},
  fonttitle=\small\bfseries]
\begin{theorem}[Unique Syntax Tree for Propositional Formulas]
\label{thm:unique-syntax-tree-propositional-formulas}
Every propositional well-formed formula has a unique parse tree.

\noindent\hyperref[prf:unique-syntax-tree-propositional-formulas]{Go to proof.}
\end{theorem}
\end{tcolorbox}

\begin{remark*}[Standard quantified statement]
\[
\forall\varphi\in\WFF,\quad \exists!T\; (T=\operatorname{Tree}(\varphi)).
\]
\end{remark*}

\begin{remark*}[Interpretation]
A formula has exactly one tree representation of its recursive construction. This is the tree form of unique decomposition.
\end{remark*}

\begin{dependencies}
\begin{itemize}
  \item \hyperref[def:parse-tree-propositional-logic]{Parse Tree}
  \item \hyperref[thm:unique-decomposition-propositional-formulas]{Unique Decomposition for Propositional Formulas}
\end{itemize}
\end{dependencies}
% END GENERATED ARTIFACT: thm:unique-syntax-tree-propositional-formulas
% BEGIN GENERATED ARTIFACT: def:formula-variable-set-propositional-logic
\begin{tcolorbox}[colback=propbox, colframe=propborder, arc=2pt,
  left=6pt, right=6pt, top=4pt, bottom=4pt,
  title={\small\textbf{Definition (Variable Set of a Formula)}},
  fonttitle=\small\bfseries]
\begin{definition}[Variable Set of a Formula]
\label{def:formula-variable-set-propositional-logic}
The variable set \(\operatorname{Var}(\varphi)\) of a formula \(\varphi\) is defined recursively:
\begin{enumerate}
\item if \(\varphi=P\in\Prop\), then \(\operatorname{Var}(\varphi)=\{P\}\);
\item if \(\varphi=\neg\psi\), then \(\operatorname{Var}(\varphi)=\operatorname{Var}(\psi)\);
\item if \(\varphi=(\psi\circ\chi)\), then
\[
\operatorname{Var}(\varphi)=\operatorname{Var}(\psi)\cup\operatorname{Var}(\chi).
\]
\end{enumerate}
\end{definition}
\end{tcolorbox}

\begin{remark*}[Standard quantified statement]
\[
\operatorname{Var}(P)=\{P\},\qquad
\operatorname{Var}(\neg\psi)=\operatorname{Var}(\psi),\qquad
\operatorname{Var}((\psi\circ\chi))=\operatorname{Var}(\psi)\cup\operatorname{Var}(\chi).
\]
\end{remark*}

\begin{remark*}[Interpretation]
The variable set records exactly which propositional variables occur in a formula.
\end{remark*}

\begin{dependencies}
\begin{itemize}
  \item \hyperref[def:well-formed-formula]{Well-Formed Formula}
  \item \hyperref[thm:structural-recursion-propositional-formulas]{Structural Recursion for Propositional Formulas}
\end{itemize}
\end{dependencies}
% END GENERATED ARTIFACT: def:formula-variable-set-propositional-logic
% BEGIN GENERATED ARTIFACT: def:subformula-propositional-logic
\begin{tcolorbox}[colback=propbox, colframe=propborder, arc=2pt,
  left=6pt, right=6pt, top=4pt, bottom=4pt,
  title={\small\textbf{Definition (Subformula)}},
  fonttitle=\small\bfseries]
\begin{definition}[Subformula]
\label{def:subformula-propositional-logic}
The set \(\operatorname{Sub}(\varphi)\) of subformulas of a formula \(\varphi\) is defined recursively:
\begin{enumerate}
\item if \(\varphi\) is atomic, then \(\operatorname{Sub}(\varphi)=\{\varphi\}\);
\item if \(\varphi=\neg\psi\), then
\[
\operatorname{Sub}(\varphi)=\{\varphi\}\cup\operatorname{Sub}(\psi);
\]
\item if \(\varphi=(\psi\circ\chi)\), then
\[
\operatorname{Sub}(\varphi)=\{\varphi\}\cup\operatorname{Sub}(\psi)\cup\operatorname{Sub}(\chi).
\]
\end{enumerate}
\end{definition}
\end{tcolorbox}

\begin{remark*}[Standard quantified statement]
\[
\operatorname{Sub}(P)=\{P\},\qquad
\operatorname{Sub}(\neg\psi)=\{\neg\psi\}\cup\operatorname{Sub}(\psi),
\]
\[
\operatorname{Sub}((\psi\circ\chi))
=
\{(\psi\circ\chi)\}\cup\operatorname{Sub}(\psi)\cup\operatorname{Sub}(\chi).
\]
\end{remark*}

\begin{tcolorbox}[colback=propbox, colframe=propborder, arc=2pt,
  left=6pt, right=6pt, top=4pt, bottom=4pt,
  title={\small\textbf{Definition (Proper Subformula)}},
  fonttitle=\small\bfseries]
\begin{definition}[Proper Subformula]
\label{def:proper-subformula}
Let \(\varphi\in\WFF\). A \emph{proper subformula} of \(\varphi\) is a
formula \(\psi\) such that
\[
\psi\in\operatorname{Sub}(\varphi)
\quad\text{and}\quad
\psi\neq\varphi.
\]
\end{definition}
\end{tcolorbox}

\begin{remark*}[Standard quantified statement]
\[
\operatorname{ProperSubformula}(\psi,\varphi)
\Longleftrightarrow
\psi\in\operatorname{Sub}(\varphi)\land\psi\neq\varphi.
\]
\end{remark*}

\begin{remark*}[Interpretation]
Proper subformulas are the strict syntactic parts of a formula. They exclude the
whole formula itself.
\end{remark*}

\begin{dependencies}
\begin{itemize}
  \item \hyperref[def:subformula-propositional-logic]{Subformula}
  \item \hyperref[def:well-formed-formula]{Well-Formed Formula}
\end{itemize}
\end{dependencies}

\begin{remark*}[Interpretation]
Subformulas are the formulas that occur at nodes of the parse tree.
\end{remark*}

\begin{dependencies}
\begin{itemize}
  \item \hyperref[def:well-formed-formula]{Well-Formed Formula}
  \item \hyperref[thm:unique-decomposition-propositional-formulas]{Unique Decomposition for Propositional Formulas}
\end{itemize}
\end{dependencies}
% END GENERATED ARTIFACT: def:subformula-propositional-logic
% BEGIN GENERATED ARTIFACT: def:formula-depth-propositional-logic
\begin{tcolorbox}[colback=propbox, colframe=propborder, arc=2pt,
  left=6pt, right=6pt, top=4pt, bottom=4pt,
  title={\small\textbf{Definition (Formula Depth)}},
  fonttitle=\small\bfseries]
\begin{definition}[Formula Depth]
\label{def:formula-depth-propositional-logic}
The depth of a formula \(\varphi\), denoted \(\operatorname{depth}(\varphi)\), is defined recursively:
\begin{enumerate}
\item if \(\varphi\) is atomic, then \(\operatorname{depth}(\varphi)=0\);
\item if \(\varphi=\neg\psi\), then \(\operatorname{depth}(\varphi)=\operatorname{depth}(\psi)+1\);
\item if \(\varphi=(\psi\circ\chi)\), then
\[
\operatorname{depth}(\varphi)=\max\{\operatorname{depth}(\psi),\operatorname{depth}(\chi)\}+1.
\]
\end{enumerate}
\end{definition}
\end{tcolorbox}

\begin{remark*}[Standard quantified statement]
\[
\operatorname{depth}(P)=0,\qquad
\operatorname{depth}(\neg\psi)=\operatorname{depth}(\psi)+1,
\]
\[
\operatorname{depth}((\psi\circ\chi))
=
\max\{\operatorname{depth}(\psi),\operatorname{depth}(\chi)\}+1.
\]
\end{remark*}

\begin{remark*}[Interpretation]
Formula depth measures the height of the syntax tree. Some authors call this the rank, complexity, or formation depth of the formula.
\end{remark*}

\begin{dependencies}
\begin{itemize}
  \item \hyperref[def:well-formed-formula]{Well-Formed Formula}
  \item \hyperref[def:parse-tree-propositional-logic]{Parse Tree}
\end{itemize}
\end{dependencies}
% END GENERATED ARTIFACT: def:formula-depth-propositional-logic
% BEGIN GENERATED ARTIFACT: def:connective-count-propositional-logic
\begin{tcolorbox}[colback=propbox, colframe=propborder, arc=2pt,
  left=6pt, right=6pt, top=4pt, bottom=4pt,
  title={\small\textbf{Definition (Connective Count)}},
  fonttitle=\small\bfseries]
\begin{definition}[Connective Count]
\label{def:connective-count-propositional-logic}
The connective count of a formula \(\varphi\), denoted \(\operatorname{conn}(\varphi)\), is defined recursively:
\begin{enumerate}
\item if \(\varphi\) is atomic, then \(\operatorname{conn}(\varphi)=0\);
\item if \(\varphi=\neg\psi\), then \(\operatorname{conn}(\varphi)=\operatorname{conn}(\psi)+1\);
\item if \(\varphi=(\psi\circ\chi)\), then
\[
\operatorname{conn}(\varphi)=\operatorname{conn}(\psi)+\operatorname{conn}(\chi)+1.
\]
\end{enumerate}
\end{definition}
\end{tcolorbox}

\begin{remark*}[Standard quantified statement]
\[
\operatorname{conn}(P)=0,\qquad
\operatorname{conn}(\neg\psi)=\operatorname{conn}(\psi)+1,
\]
\[
\operatorname{conn}((\psi\circ\chi))
=
\operatorname{conn}(\psi)+\operatorname{conn}(\chi)+1.
\]
\end{remark*}

\begin{remark*}[Interpretation]
Formula depth measures the height of the syntax tree. Connective count measures the number of connective nodes in the syntax tree.
\end{remark*}

\begin{dependencies}
\begin{itemize}
  \item \hyperref[def:well-formed-formula]{Well-Formed Formula}
  \item \hyperref[def:parse-tree-propositional-logic]{Parse Tree}
\end{itemize}
\end{dependencies}
% END GENERATED ARTIFACT: def:connective-count-propositional-logic
% BEGIN GENERATED ARTIFACT: lem:proper-subformula-depth-propositional-logic
\begin{tcolorbox}[colback=lembox, colframe=lemborder, arc=2pt,
  left=6pt, right=6pt, top=4pt, bottom=4pt,
  title={\small\textbf{Lemma (Proper Subformula Depth)}},
  fonttitle=\small\bfseries]
\begin{lemma}[Proper Subformula Depth]
\label{lem:proper-subformula-depth-propositional-logic}
If \(\psi\) is a proper subformula of \(\varphi\), then
\[
\operatorname{depth}(\psi)<\operatorname{depth}(\varphi).
\]

\noindent\hyperref[prf:proper-subformula-depth-propositional-logic]{Go to proof.}
\end{lemma}
\end{tcolorbox}

\begin{remark*}[Standard quantified statement]
\[
(\forall\varphi,\psi\in\WFF)
\left[
\psi\in\operatorname{Sub}(\varphi)\land \psi\ne\varphi
\to
\operatorname{depth}(\psi)<\operatorname{depth}(\varphi)
\right].
\]
\end{remark*}

\begin{remark*}[Interpretation]
Every descent into a proper subformula lowers formula depth. This justifies induction on formula depth and recursive proofs over subformulas.
\end{remark*}

\begin{dependencies}
\begin{itemize}
  \item \hyperref[def:subformula-propositional-logic]{Subformula}
  \item \hyperref[def:formula-depth-propositional-logic]{Formula Depth}
\end{itemize}
\end{dependencies}
% END GENERATED ARTIFACT: lem:proper-subformula-depth-propositional-logic
% BEGIN GENERATED ARTIFACT: lem:finiteness-of-variables-propositional-formulas
\begin{tcolorbox}[colback=lembox, colframe=lemborder, arc=2pt,
  left=6pt, right=6pt, top=4pt, bottom=4pt,
  title={\small\textbf{Lemma (Finiteness of Variables in a Formula)}},
  fonttitle=\small\bfseries]
\begin{lemma}[Finiteness of Variables in a Formula]
\label{lem:finiteness-of-variables-propositional-formulas}
For every \(\varphi\in\WFF\), the set \(\operatorname{Var}(\varphi)\) is finite.

\noindent\hyperref[prf:finiteness-of-variables-propositional-formulas]{Go to proof.}
\end{lemma}
\end{tcolorbox}

\begin{remark*}[Standard quantified statement]
\[
(\forall\varphi\in\WFF)\quad
\operatorname{Var}(\varphi)\text{ is finite}.
\]
\end{remark*}

\begin{remark*}[Interpretation]
Every formula depends on only finitely many propositional variables, even though the language may contain infinitely many variables.
\end{remark*}

\begin{dependencies}
\begin{itemize}
  \item \hyperref[def:formula-variable-set-propositional-logic]{Variable Set of a Formula}
  \item \hyperref[thm:structural-induction-propositional-formulas]{Structural Induction for Propositional Formulas}
\end{itemize}
\end{dependencies}
% END GENERATED ARTIFACT: lem:finiteness-of-variables-propositional-formulas
% BEGIN GENERATED ARTIFACT: lem:finiteness-of-subformulas-propositional-formulas
\begin{tcolorbox}[colback=lembox, colframe=lemborder, arc=2pt,
  left=6pt, right=6pt, top=4pt, bottom=4pt,
  title={\small\textbf{Lemma (Finiteness of Subformulas)}},
  fonttitle=\small\bfseries]
\begin{lemma}[Finiteness of Subformulas]
\label{lem:finiteness-of-subformulas-propositional-formulas}
For every \(\varphi\in\WFF\), the set \(\operatorname{Sub}(\varphi)\) is finite.

\noindent\hyperref[prf:finiteness-of-subformulas-propositional-formulas]{Go to proof.}
\end{lemma}
\end{tcolorbox}

\begin{remark*}[Standard quantified statement]
\[
(\forall\varphi\in\WFF)\quad
\operatorname{Sub}(\varphi)\text{ is finite}.
\]
\end{remark*}

\begin{remark*}[Interpretation]
Every formula has only finitely many syntactic parts.
\end{remark*}

\begin{dependencies}
\begin{itemize}
  \item \hyperref[def:subformula-propositional-logic]{Subformula}
  \item \hyperref[thm:structural-induction-propositional-formulas]{Structural Induction for Propositional Formulas}
\end{itemize}
\end{dependencies}
% END GENERATED ARTIFACT: lem:finiteness-of-subformulas-propositional-formulas
